% Recourse function section — paste after the main model section
% Required packages: amsmath, amssymb
 
\subsection{Recourse Function}
 
The second-stage problem is formulated as the recourse function $Q(x, s)$,
which gives the minimum operational cost achievable for a fixed first-stage
solution $x = (A_p, X_{pp'h}, Y_{ph}, Z_h, V_{hq})$ and a realised scenario
$s$ with surgery durations $d_{ps}$:
%
\begin{align}
  Q(x, s) = \min_{S,W,I,O} \quad
    & \beta_1 \sum_{p \in P} W_{ps}
    + \beta_2 \sum_{h \in H} I_{hs}
    + \beta_3 \sum_{h \in H} O_{hs}
  \label{eq:recourse_obj}
\end{align}
 
\noindent subject to:
 
\paragraph{Start time propagation.}
If patient $p'$ follows patient $p$ in session $h$, the start of $p'$ must
respect patient $p$'s realised duration and the changeover time $c$:
%
\begin{align}
  S_{p's} \geq\; & S_{ps} + d_{ps} + c - M\bigl(1 - X_{pp'h}\bigr)
  \nonumber \\
                 & \forall\; p, p' \in P,\; p \neq p',\; h \in H
  \label{eq:rc_seq}
\end{align}
 
\paragraph{First patient lower bound.}
The first patient in session $h$ starts no earlier than the session opening:
%
\begin{align}
  S_{ps} \geq\; t^{\text{start}} - M\bigl(1 - X_{0ph}\bigr)
  \qquad \forall\; p \in P,\; h \in H
  \label{eq:rc_first}
\end{align}
 
\paragraph{Appointment lower bound.}
A patient cannot start surgery before their appointed time:
%
\begin{align}
  S_{ps} \geq\; A_p - M(1 - Y_{ph})
  \qquad \forall\; p \in P,\; h \in H
  \label{eq:rc_appt}
\end{align}
 
\paragraph{Waiting time.}
Waiting time is the excess of actual start time over appointment time:
%
\begin{equation}
  W_{ps} \geq S_{ps} - A_p
  \qquad \forall\; p \in P
  \label{eq:rc_wait}
\end{equation}
 
\paragraph{Overtime.}
Overtime is incurred when any patient's finish time exceeds the session closing
time. The big-$M$ term deactivates the constraint for patients not in session
$h$:
%
\begin{align}
  O_{hs} \geq\; S_{ps} + d_{ps} - t^{\text{close}} - M(1 - Y_{ph})
  \qquad \forall\; p \in P,\; h \in H
  \label{eq:rc_ot}
\end{align}
 
\paragraph{Idle time.}
Idle time equals the session window minus total utilisation, corrected for
overtime. There are $\sum_{p}Y_{ph} - Z_h$ inter-patient changeovers in an
open session (one fewer than the number of patients):
%
\begin{align}
  I_{hs} \geq\; & \bigl(t^{\text{close}} - t^{\text{start}}\bigr) Z_h
                  - \sum_{p \in P} d_{ps}\, Y_{ph} \nonumber \\
                & - c\Bigl(\sum_{p \in P} Y_{ph} - Z_h\Bigr)
                  + O_{hs}
  \qquad \forall\; h \in H
  \label{eq:rc_idle}
\end{align}
 
\paragraph{Non-negativity.}
%
\begin{equation}
  S_{ps},\; W_{ps},\; I_{hs},\; O_{hs} \geq 0
  \qquad \forall\; p \in P,\; h \in H
  \label{eq:rc_nn}
\end{equation}
 
\noindent The first-stage variables $A_p$, $X_{pp'h}$, $Y_{ph}$, $Z_h$ enter
$Q(x,s)$ as fixed parameters. Given these, all constraints are linear in the
second-stage variables, so $Q(x, s)$ is a linear program for each $(x, s)$.
 
\subsection{Deterministic Equivalent}
 
Rather than solving $Q(x,s)$ separately for each scenario, the two-stage
model is reformulated as a single deterministic equivalent (DE) by replicating
the second-stage variables and constraints across all scenarios $s \in S$. The
full model then becomes:
%
\begin{align}
  \min \quad
    & \beta_4 \sum_{h \in H} \sum_{q \in Q} D_{hq}
    + \sum_{s \in S} \pi_s \, Q(x, s)
  \label{eq:de_obj}
\end{align}
 
\noindent which, after substituting the recourse function, expands to the
compact form in \eqref{eq:obj}, subject to all first-stage constraints
\eqref{eq:appt_bounds}--\eqref{eq:mtz} and, for each scenario $s \in S$,
the recourse constraints \eqref{eq:rc_seq}--\eqref{eq:rc_nn}.
 
The DE is a mixed-integer linear program (MILP): all integrality requirements
appear in the first stage, and the second-stage subproblem is a pure LP for
any fixed first-stage solution. The total number of second-stage variables
scales as $\mathcal{O}(|S| \cdot (|P| + |H|))$, so tractability depends on
keeping $|S|$ moderate, for instance through Sample Average Approximation
(SAA), in which $|S|$ scenarios are drawn from the duration distribution and
the DE is solved once.
